\newtcolorbox{textBox}[1][]{
 parbox=false,
 			colback=Teal!10,
 			colframe=Teal,
 			colbacktitle=Teal,
 			enhanced jigsaw,
 			breakable,  
 			attach boxed title to top text left={yshift=-3mm,yshifttext=-3mm}, 
 			#1,%
}

\section*{\emoji{input-numbers} Zusammensetzung der Note}

Sie werden im Laufe des Semesters soviele schriftliche Prüfungen ablegen, wie sie Wochenlektionen haben, also drei Prüfungen im ersten Semester. Sämtliche schriftlichen Noten werden zu gleichen Teilen gewichtet.

Bitte beachten Sie, dass wir Ihnen die Noten aller Leistungskontrollen (schriftliche Prüfungsnoten, mündliche Note) bereits auf 1/10 gerundet bekanntgeben werden  (z.B. 4.9 oder 5.3). Wir werden dann stets mit diesen auf Zehntel gerundeten Noten-Werten arbeiten.

\begin{mycaution}
Wer zu einem festgesetzten Probentermin nicht erscheint, schreibt am Semesterende eine \textbf{Semesterprobe}. Es erfolgt keine spezielle Einladung, da der Termin jeweils zu Beginn des Semesters kommuniziert wird. Es liegt im Ermessen der Lehrperson, ob in begründeten Fällen (z.B. amtliches Aufgebot, länger dauernde Krankheit, o.Ä.) eine Ausnahme von dieser Regelung gewährt werden kann. Das Datum der nächsten Semesterprobe ist: \textbf{Freitag, 20.06.2025, Zimmer 401}.
\end{mycaution}

\begin{textBox}[title=\emoji{speaking-head} Bonus 1: Wortmeldungen im Unterricht]

Die mündliche Note ergibt sich aus der (subjektiv wahrgenommenen) Qualität und Quantität (Menge) Ihrer Wortmeldungen im Unterricht. Der Mündlich-Bonus $B_m$ wird wie folgt berechnet:

\begin{table}[H]
\centering
\begin{tabular}{c|c}
\textbf{Mündliches Engagement} & \textbf{Noten-Bonus} \\\midrule\midrule
Sehr hoch & 0.15 \\\midrule
Hoch & 0.1 \\\midrule
Durchschnittlich & 0 \\
\end{tabular}
\end{table}

\end{textBox}

\begin{textBox}[title=\emoji{speech-balloon} Bonus 2: Rückmeldungen zu den Skripts und Unterlagen]
Aktuell befinden sich viele meiner Skripts im Aufbau, daher können sich an verschieden Orten Fehler eingeschlichen haben. Sie haben die Möglichkeit, Ihre \textit{Endnote} um 0.05 Punkte zu erhöhen, indem Sie insgesamt \textbf{10 Punkte} durch Rückmeldungen sammeln. 0.05 Notenpunkte klingt nach wenig, dies führt jedoch in 20\% aller Fälle zu einem Anstieg einer halben Note \textit{im Zeugnis}!

Sie können pro akzeptierte Einreichung maximal 2 Punkte sammeln:
\begin{itemize}
\item \textbf{1 Punkt} für Schreibfehler, Tippfehler, stilistisch unsaubere Sätze
\item \textbf{2 Punkte} für inhaltliche Fehler oder für inhaltliche Vorschläge, wie das Skript verbessert werden kann (z.B. Abschnitte, die man verkürzen kann oder Dinge, die man noch hinzufügen sollte)
\end{itemize}

Rückmeldungen zu Moodle-Seiten oder anderen, von mir erstellten Unterlagen zählen ebenfalls! Jede Rückmeldung wird nur einmal beachtet, es gilt also das Prinzip ``der / die Schnellere erhält die Punkte''.

%Um meinen Arbeitsaufwand überschaubar zu halten, vermeiden Sie bitte das Einreichen von Meldungen, bei denen Sie nicht 100\% sicher sind. Insbesondere bei Punkt 3 sollten Sie nur Dinge vorschlagen, die wirklich zentral sind. 

\textbf{Wichtig}: Damit Sie Punkte erhalten, müssen Sie immer klar angeben, wo Sie den Fehler gefunden haben. Reichen Sie Ihre Meldungen stets über das \href{}{Rückmeldungs-Formular auf Moodle} ein!

Darüber, welche Vorschläge umgesetzt werden und welche Punkte angerechnet werden, entscheidet die Lehrperson. Ich danke Ihnen im Voraus für Ihre Rückmeldungen!



\end{textBox}

Die finale Semester-Endnote $N_e$ (Note, die im Zeugnis stehen wird) ergibt sich in diesem Fall aus dem schriftlichen Schnitt $N_s$ sowie dem Mündlich-Bonus $B_m$ und dem Feedback-Bonus $B_f$:

\[N_e=\text{\texttt{runden\_auf\_halbe\_note}}\Big(N_s+B_m+B_f\Big).\]